\documentclass[aspectratio=169,11pt]{beamer}

% ---------- Theme and appearance ----------
\usetheme{Madrid}
\usecolortheme{default}
\setbeamertemplate{navigation symbols}{}
\setbeamertemplate{caption}[numbered]
\setbeamertemplate{itemize items}[circle]
\setbeamertemplate{enumerate items}[default]
\setbeamercolor{title}{fg=white,bg=blue!70!black}
\setbeamercolor{frametitle}{fg=white,bg=blue!70!black}
\setbeamercolor{structure}{fg=blue!60!black}
\setbeamercolor{block title}{fg=white,bg=blue!70!black}
\setbeamercolor{block body}{bg=blue!5}
\setbeamertemplate{blocks}[rounded][shadow=false]

% ---------- Packages ----------
\usepackage[utf8]{inputenc}
\usepackage[T1]{fontenc}
\usepackage{lmodern}
\usepackage{listings}
\usepackage{xcolor}
\usepackage{hyperref}
\usepackage{graphicx}
\usepackage{booktabs}
\usepackage{array}
\usepackage{amsmath,amssymb}

\lstset{
	language=Python,
	basicstyle=\ttfamily\scriptsize,
	breaklines=true,
	columns=fullflexible
}
% ---------- Code listing style ----------
\definecolor{codebg}{RGB}{245,245,245}
\definecolor{codecomment}{RGB}{0,128,0}
\definecolor{codekeyword}{RGB}{0,0,255}
\definecolor{codestring}{RGB}{163,21,21}
\lstdefinestyle{pythonstyle}{
	language=Python,
	backgroundcolor=\color{codebg},
	basicstyle=\ttfamily\footnotesize,
	keywordstyle=\color{codekeyword}\bfseries,
	commentstyle=\color{codecomment}\itshape,
	stringstyle=\color{codestring},
	numbers=left,
	numberstyle=\tiny\color{gray},
	numbersep=6pt,
	frame=single,
	rulecolor=\color{gray!40},
	breaklines=true,
	showstringspaces=false,
	tabsize=4,
	captionpos=b
}
\lstset{style=pythonstyle}

% ---------- Title metadata ----------
\title[Notebook Cleanup \& Teradata Analysis]{Demystifying a Legacy Teradata Notebook}
\subtitle{Breaking a Monolithic Jupyter Notebook into Clean, Readable Cells}
\author{Your Name}
\institute{Your Department / Course}
\date{\today}

% ---------- Document ----------
\begin{document}
	
	% ---------- Title frame ----------
	\begin{frame}
		\titlepage
	\end{frame}
	
	% ---------- Outline ----------
	\begin{frame}{Outline}
		\tableofcontents
	\end{frame}
	
	% ============================================================
	\section{Introduction and Context}
	% ============================================================
	
	\begin{frame}{What This Project Is About}
		\begin{itemize}
			\item We are given a \textbf{Jupyter Notebook} that begins with a \textbf{Teradata connection}.
			\item The notebook lives on a \textbf{Jupyter server}.
			\item The notebook is \textbf{not} a clean, modular notebook.
			\item It contains \textbf{one very large cell} with:
			\begin{itemize}
				\item A large SQL query,
				\item Python \texttt{for} loops that inject strings into the SQL,
				\item Multiple smaller \texttt{pandas} DataFrames,
				\item Business logic mixed with data access.
			\end{itemize}
			\item Goal: \textbf{demystify}, \textbf{break down}, and \textbf{clean up} the notebook.
		\end{itemize}
	\end{frame}
	
	\begin{frame}{Why This Matters}
		\begin{block}{The Problem}
			A single monolithic cell is hard to:
			\begin{itemize}
				\item Read,
				\item Debug,
				\item Explain,
				\item Reuse,
				\item Hand over to others (or to Copilot).
			\end{itemize}
		\end{block}
		
		\begin{block}{The Solution}
			Split the code into smaller, logical cells and save intermediate outputs (e.g., DataFrame headers) so the workflow becomes transparent.
		\end{block}
	\end{frame}
	
	\begin{frame}{What We Are \emph{Not} Doing}
		\begin{itemize}
			\item We are \textbf{not} rewriting the business logic.
			\item We are \textbf{not} changing the SQL semantics.
			\item We are \textbf{not} using Copilot for this phase.
		\end{itemize}
		
		\begin{alertblock}{Why avoid Copilot for now?}
			We are preparing clean, structured input \emph{for} Copilot.\\
			Using Copilot too early adds noise and confusion.
		\end{alertblock}
	\end{frame}
	
	% ============================================================
	\section{Environment and Access}
	% ============================================================
	
	\begin{frame}{Standard Procedure to Run the Notebook}
		\begin{enumerate}
			\item Log in to the server using \textbf{PuTTY}.
			\item Switch to the \textbf{Power Broker} account.
			\item Navigate to our project folder.
			\item Run the notebook command.
			\item The notebook fetches the password via \texttt{config.py}.
			\item Connect to \textbf{Teradata}.
		\end{enumerate}
		
		\begin{block}{Note}
			Teradata is a legacy, classical data warehouse with its own computational framework.\\
			For this course, it is important to understand how it differs from \textbf{Hive}.
		\end{block}
	\end{frame}
	
	\begin{frame}[fragile]{Typical Access Flow}
		\begin{lstlisting}[language=bash,caption={Shell commands (illustrative)}]
			# Log in via PuTTY (outside the notebook)
			# Then on the server:
			sudo su - powerbroker
			cd /path/to/our/folder
			jupyter nbconvert --to notebook --execute our_notebook.ipynb
		\end{lstlisting}
	\end{frame}
	
	% ============================================================
	\section{The Monolithic Notebook}
	% ============================================================
	
	\begin{frame}[fragile]{What the Big Cell Looks Like}
		\begin{lstlisting}[caption={Conceptual structure of the monolithic cell}]
			# config.py provides the password
			from config import get_password
			
			password = get_password()
			conn = teradata.connect(..., password=password)
			
			# One huge query with Python string injection
			base_query = """
			SELECT ...
			FROM ...
			WHERE month = '{month}'
			"""
			
			for month in months:
			query = base_query.format(month=month)
			df = pd.read_sql(query, conn)
			# ... lots of pandas logic ...
			# ... more smaller DataFrames ...
			# ... final output ...
		\end{lstlisting}
	\end{frame}
	
	\begin{frame}{Why the Query Is So Big}
		\begin{itemize}
			\item To \textbf{iteratively inject strings} into SQL, we need Python.
			\item A \texttt{for} loop (e.g., over months) injects a value into the query.
			\item Each iteration produces a separate query result.
			\item The original author wrote the loop and the query together.
			\item Over time, more and more logic was added into the same cell.
		\end{itemize}
		
		\begin{block}{Key Insight}
			The big query is not big because the SQL is complex.\\
			It is big because \textbf{Python and SQL are entangled}.
		\end{block}
	\end{frame}
	
	\begin{frame}{What Else Is Inside the Big Cell?}
		\begin{itemize}
			\item Multiple \texttt{pandas} DataFrames are created.
			\item Each DataFrame is separated into its own logic.
			\item Intermediate results are used to build the next step.
			\item The final output is saved or returned.
		\end{itemize}
		
		\begin{alertblock}{Our Job}
			Separate the logic into different blocks and find out what is happening at each step.
		\end{alertblock}
	\end{frame}
	
	% ============================================================
	\section{Cleanup Strategy}
	% ============================================================
	
	\begin{frame}{Step-by-Step Cleanup Plan}
		\begin{enumerate}
			\item \textbf{Identify} the major logical sections in the big cell.
			\item \textbf{Split} them into smaller Jupyter cells.
			\item \textbf{Preserve} the original code (do not change logic).
			\item \textbf{Save} intermediate outputs:
			\begin{itemize}
				\item DataFrame headers,
				\item Shapes,
				\item Sample rows,
				\item Column definitions.
			\end{itemize}
			\item \textbf{Document} what each cell does.
		\end{enumerate}
	\end{frame}
	
	\begin{frame}{Example: Splitting the Big Cell}
		\begin{columns}[T]
			\begin{column}{0.48\textwidth}
				\textbf{Before}
				\begin{itemize}
					\item One cell
					\item 500+ lines
					\item SQL + Python + pandas
					\item Hard to debug
				\end{itemize}
			\end{column}
			\begin{column}{0.48\textwidth}
				\textbf{After}
				\begin{itemize}
					\item Cell 1: Imports and config
					\item Cell 2: Connection
					\item Cell 3: Query template
					\item Cell 4: Loop over months
					\item Cell 5: DataFrame 1
					\item Cell 6: DataFrame 2
					\item Cell 7: Final output
				\end{itemize}
			\end{column}
		\end{columns}
	\end{frame}
	
	\begin{frame}[fragile]{Example: Cleaned Cell Structure}
		\begin{lstlisting}[caption={Split cells (conceptual)}]
			# Cell 1: Imports
			import pandas as pd
			from config import get_password
			
			# Cell 2: Connection
			password = get_password()
			conn = teradata.connect(..., password=password)
			
			# Cell 3: Query template
			base_query = """
			SELECT ...
			FROM ...
			WHERE month = '{month}'
			"""
			
			# Cell 4: Loop over months
			monthly_dfs = {}
			for month in months:
			query = base_query.format(month=month)
			monthly_dfs[month] = pd.read_sql(query, conn)
			
			# Cell 5: Inspect one DataFrame
			df = monthly_dfs[months[0]]
			print(df.head())
			print(df.columns.tolist())
		\end{lstlisting}
	\end{frame}
	
	\begin{frame}{Saving Intermediate Outputs}
		\begin{itemize}
			\item For each DataFrame, save:
			\begin{itemize}
				\item \texttt{df.head()},
				\item \texttt{df.shape},
				\item \texttt{df.columns.tolist()},
				\item \texttt{df.dtypes}.
			\end{itemize}
			\item This makes it easy to:
			\begin{itemize}
				\item Verify data,
				\item Compare months,
				\item Explain to others,
				\item Prepare input for Copilot.
			\end{itemize}
		\end{itemize}
		
		\begin{block}{Remember}
			We are \textbf{not} touching the code. We are only \textbf{reorganizing} it.
		\end{block}
	\end{frame}
	
\begin{frame}{Defining Columns at the Start}
	
	\begin{itemize}
		\item At the top of the notebook, define the columns we care about.
		\item Example:
	\end{itemize}
	
	\begin{block}{Column definitions}
		\small
		\texttt{\detokenize{KEY_COLUMNS = [}}
		\par
		\texttt{\detokenize{    "customer_id",}}
		\par
		\texttt{\detokenize{    "month",}}
		\par
		\texttt{\detokenize{    "revenue",}}
		\par
		\texttt{\detokenize{    "product_category",}}
		\par
		\texttt{\detokenize{]}}
		\par
		\vspace{0.1cm}
		\texttt{\detokenize{DATE_COLUMNS = ["month", "transaction_date"]}}
	\end{block}
	
	\begin{itemize}
		\item This makes the notebook self-documenting.
		\item It also helps when splitting the big cell.
	\end{itemize}
	
\end{frame}
	
	% ============================================================
	\section{Teradata vs. Hive}
	% ============================================================
	
	\begin{frame}{Teradata vs. Hive: Why It Matters}
		\begin{table}
			\centering
			\begin{tabular}{p{0.28\textwidth} p{0.28\textwidth} p{0.28\textwidth}}
				\toprule
				\textbf{Aspect} & \textbf{Teradata} & \textbf{Hive} \\
				\midrule
				Type & Classical RDBMS / MPP & Hadoop-based data warehouse \\
				Language & SQL (Teradata dialect) & HiveQL (SQL-like) \\
				Compute & Proprietary engine & MapReduce / Tez / Spark \\
				Storage & Proprietary & HDFS / cloud storage \\
				Use case & Legacy enterprise DW & Big data analytics \\
				\bottomrule
			\end{tabular}
		\end{table}
		
		\begin{block}{Takeaway}
			Understanding both helps you know \emph{why} the code is written the way it is.
		\end{block}
	\end{frame}
	
	\begin{frame}{Why We Still Use Teradata}
		\begin{itemize}
			\item It is a \textbf{very old, classical, legacy} place to keep data.
			\item It has its own \textbf{computational framework}.
			\item Many enterprises still rely on it for critical reporting.
			\item The notebook we are cleaning is part of that legacy pipeline.
		\end{itemize}
		
		\begin{alertblock}{Be Careful}
			Teradata SQL dialect may differ from Hive SQL.\\
			Do not assume a query will run the same way on both.
		\end{alertblock}
	\end{frame}
	
	% ============================================================
	\section{Practical Workflow}
	% ============================================================
	
	\begin{frame}{Recommended Workflow}
		\begin{enumerate}
			\item Log in via PuTTY.
			\item Switch to Power Broker account.
			\item Navigate to the project folder.
			\item Open the notebook.
			\item Create a \textbf{copy} of the original notebook.
			\item In the copy, split the big cell into smaller cells.
			\item Run each cell and save intermediate outputs.
			\item Add markdown explanations for each section.
			\item Compare the final output with the original.
		\end{enumerate}
	\end{frame}
	
	\begin{frame}[fragile]{Example: Saving Outputs to Disk}
		\begin{lstlisting}[caption={Save DataFrame headers and samples}]
			# After splitting, for each important DataFrame:
			df.head(20).to_csv("output/df_head.csv", index=False)
			with open("output/df_columns.txt", "w") as f:
			f.write("\n".join(df.columns.tolist()))
			
			# Also save a small summary
			with open("output/df_summary.txt", "w") as f:
			f.write(f"Shape: {df.shape}\n")
			f.write(f"Columns: {df.columns.tolist()}\n")
			f.write(f"Dtypes:\n{df.dtypes}\n")
		\end{lstlisting}
	\end{frame}
	
	\begin{frame}{What Good Looks Like}
		\begin{itemize}
			\item A notebook with \textbf{clear sections}:
			\begin{itemize}
				\item Imports,
				\item Configuration,
				\item Connection,
				\item Query templates,
				\item Loops,
				\item DataFrame processing,
				\item Final output.
			\end{itemize}
			\item Each cell does \textbf{one thing}.
			\item Intermediate outputs are \textbf{saved and visible}.
			\item The original logic is \textbf{unchanged}.
		\end{itemize}
	\end{frame}
	
	% ============================================================
	\section{Summary and Next Steps}
	% ============================================================
	
	\begin{frame}{Summary}
		\begin{itemize}
			\item We start with a monolithic Jupyter notebook connected to Teradata.
			\item The big cell mixes SQL, Python loops, and pandas logic.
			\item Our job is to \textbf{demystify} and \textbf{clean up} the notebook.
			\item We split the big cell into smaller, logical cells.
			\item We save intermediate outputs (headers, shapes, samples).
			\item We avoid Copilot for now to keep the analysis clean.
			\item We understand Teradata vs. Hive to better interpret the code.
		\end{itemize}
	\end{frame}
	
	\begin{frame}{Next Steps}
		\begin{enumerate}
			\item Get access to the actual notebook.
			\item Identify the major logical blocks.
			\item Create a cleaned copy.
			\item Split the big cell step by step.
			\item Save outputs for each DataFrame.
			\item Document everything in markdown.
			\item Only then consider using Copilot for further analysis.
		\end{enumerate}
		
		\begin{block}{Remember}
			Clean code is not about rewriting.\\
			It is about making the existing code \textbf{understandable}.
		\end{block}
	\end{frame}
	
	\begin{frame}{Questions?}
		\begin{center}
			\Large
			Thank you!\\[1em]
			\normalsize
			Feel free to ask questions or share your notebook for a live walkthrough.
		\end{center}
	\end{frame}
	
\end{document}